\section{一个更大的直线类产生七个维数的构造}

\href{https://www.math.rwth-aachen.de/~Gabriele.Nebe/LATTICES/P48p.html}{格目录中的
\(P_{48p}\)} 是 48 维偶幺模格，最小平方范数为六。其 theta
级数确定最小壳层含 52,416,000 个向量\cite{p48p-catalogue,nebe}。

\subsection{由重合的像得到不交的块}

新的有限对象是包含 7,077 条最短直线的母类。在目录 Gram 矩阵 \(G\)
的整数基下，代表行满足

\[
cGc^t=6,\qquad |cGc'^t|\le2\quad(c\ne\pm c').
\]

利用精确格自同构得到更多直线类。选择若干个像，对重复直线按首次出现分配，得到互不相交的子集，总并集大小记为
\(S\)。每条直线取正负两个方向，便得到可用于提升的母块。

依次使用根码
\[
A_1,\ A_2,\ D_3,\ D_4,\ D_5,\ E_6,\ E_7
\]
作为一至七维尾码，其点数依次为
\[
2,\ 6,\ 12,\ 24,\ 40,\ 72,\ 126.
\]
每个母块分配到不同的反足尾直线。每条使用过的母直线删除两个赤道点，产生四个帽点；帽点的头部平方范数为
\(3/4\)，尾部平方范数为 \(1/4\)。因此

\[
K(48+k)\ge52416000-2S+4S+\tau_k=52416000+2S+\tau_k.
\]

\begin{table}[htbp]
\centering
\begin{tabular}{@{}rrrr@{}}
\toprule
\(k\) & 选取的自同构像数 & 并集大小 \(S\) & 尾点数
\(\tau_k\)\tabularnewline
\midrule
1 & 1 & 7,077 & 2\tabularnewline
2 & 3 & 21,231 & 6\tabularnewline
3 & 6 & 42,459 & 12\tabularnewline
4 & 12 & 84,874 & 24\tabularnewline
5 & 20 & 141,328 & 40\tabularnewline
6 & 36 & 253,851 & 72\tabularnewline
7 & 63 & 442,507 & 126\tabularnewline
\bottomrule
\end{tabular}
\end{table}

\subsection{全部角距约束为何保持}

将选中的母向量除以 $\sqrt6$，记为 $x$；其对应的单位尾直线代表记为 $p_i$。
替换 $\{\pm x\}$ 的四个点为
\[
 \left(\varepsilon\frac{\sqrt3}{2}x,
       \delta\frac12p_i\right),\qquad \varepsilon,\delta\in\{\pm1\}.
\]
这些点的范数均为一。同一块中不同头直线的绝对内积至多为 $1/3$；
整个极小壳中相应上界为 $1/2$。由此逐类得到帽点间的内积上界：
\[
\begin{array}{ll}
\text{同块的不同直线：}&(3/4)(1/3)+1/4=1/2,\\
\text{同头、相反尾：}&3/4-1/4=1/2,\\
\text{相反的头：}&-3/4+1/4=-1/2,\\
\text{不同块：}&(3/4)(1/2)+(1/4)(1/2)=1/2.
\end{array}
\]
最后一行使用了不同的尾直线；任意两个块共用同一尾直线时，上界会变为 $5/8$。
帽点与保留赤道的内积至多为 $\sqrt3/4$，与纯尾点的内积至多为 $1/2$。
保留赤道与纯尾码本身已满足接吻条件。因此，直线类的有限内积条件与母格最小范数
共同覆盖全部点对，并建立上述计数公式。

\subsection{区分两部分增益}

相对于公开的7069条母直线\cite{kissingnumbers}，改进来自母类的八条新增直线和
自同构像的重新选择。将选定的像作用于原类，记并集大小为 $S_0$；比较构造的并集
大小记为 $S_{\rm ref}$。则
\[
 N-N_{\rm ref}=2(S-S_0)+2(S_0-S_{\rm ref}).
\]
第一项固定像族，衡量母类扩充；第二项固定原类，衡量像族选择。

\begin{center}
\begin{tabular}{rrrr}
\toprule
维数&母类扩充&像族选择&总增益\\\midrule
49&16&0&16\\
50&48&2&50\\
51&96&18&114\\
52&192&64&256\\
53&320&154&474\\
54&572&296&868\\
55&996&414&1410\\\bottomrule
\end{tabular}
\end{center}

50维的三个像完全不交。55维中，重合使 $8\cdot63$ 个可能的新像只贡献498条不同直线，
像族选择则在原类上再贡献207条。因此，扩大单个类与扩大像的并集是相互补充的两个问题。
